%% kapa: exam=1 round=22 year=2011 period=1 form=A subject=civil_law subject_ko=민법 seq=25 number=25
%% source: Q-Net 감정평가사 기출문제 (cst00309) artlSeq=5036790 file=extracted/1008/5036790/2011년 제22회 감정평가사 1차시험 기출문제/1교시.hwp
%% conversion: kordoc 4.24.51 -> tools/gen1.py (auto); review=proofread-sonnet; flags=none
\begin{question}{25}
법률행위의 목적에 관한 설명으로 옳지 않은 것은? (다툼이 있으면 판례에 의함)
\begin{choices}
  \item 관계 법령에서 정한 한도를 초과하는 부동산 중개수수료 약정은 그 한도를 초과하는 범위에서 무효이다.
  \item 강제집행을 면할 목적으로 부동산에 허위의 근저당권설정등기를 경료하는 행위는 반사회적 법률행위로서 무효이다.
  \item 보험계약자가 다수의 보험계약을 통하여 보험금을 부정 취득할 목적으로 보험계약을 체결한 경우, 이러한 보험계약은 반사회적 법률행위로서 무효이다.
  \item 소송에서 사실대로 증언하여 줄 것을 조건으로 어떠한 급부를 할 것을 약정한 경우, 그 급부가 통상적으로 용인될 수 있는 수준을 초과한다면 무효이다.
  \item 법률행위의 성립과정에서 강박이라는 불법적 방법이 사용된 데 불과한 때에는 반사회적 법률행위로서 무효라고 할 수 없다.
\end{choices}
\end{question}
\answer{2}
